\section{The Three-Level Audit Chain}
\label{sec:audit-chain}

\bisect{} implements a three-level SHA-256 audit chain that provides
layered protection against data manipulation.

\subsection{Level 1: Census Bureau Download Hash}

\texttt{bisect fetch --year 2020 --workers 8} downloads all required
Census files and computes SHA-256 hashes immediately upon receipt.
The hashes are stored in the manifest alongside the Census Bureau's
published checksums.

\begin{verbatim}
bisect fetch --year 2020 --state NC
[INFO] Downloading pl94171_nc_2020.pl
[INFO] SHA256: a7f3bc91...
[INFO] Comparing against Census Bureau checksum: MATCH
[INFO] Downloading tl_2020_37_tract.zip
[INFO] SHA256: d4e82ca5...
[INFO] ACS B02001 5-year: 2016-2020
[INFO] SHA256: f1c0a9b7...
[INFO] Manifest written: nc_2020_fetch.manifest.json
\end{verbatim}

A party wishing to independently verify the download can fetch the
same files directly from the Census Bureau website and compare hashes.

\subsection{Level 2: Adjacency Construction Hash}

After downloading shapefiles, \bisect{} constructs the census tract
adjacency graph.
The graph construction is deterministic: given the same TIGER/Line shapefile,
the same adjacency graph is always produced.
The SHA-256 hash of the adjacency graph is stored in the manifest.

\begin{verbatim}
bisect build nc_2020 --year 2020
[INFO] Building adjacency graph for NC
[INFO] Tracts: 2,672  Edges: 7,841
[INFO] Adjacency SHA256: 8b2a7d43...
[INFO] Adjacency hash stored in manifest
\end{verbatim}

An independent party can verify the adjacency hash by running the same
adjacency construction algorithm on their independently downloaded shapefile.
Any discrepancy indicates shapefile manipulation.

\subsection{Level 3: Final Plan Hash}

The completed district assignment (one GEOID to district mapping per state)
is hashed and stored in the manifest.

\begin{verbatim}
bisect label-verify nc_2020 --year 2020
[INFO] Replaying computation from manifest
[INFO] Census download hash: MATCH
[INFO] Adjacency hash: MATCH
[INFO] Plan hash: MATCH
[INFO] Verification: PASS
\end{verbatim}

The three-level chain provides end-to-end provenance:
any modification at any level propagates into a hash mismatch at
that level and all subsequent levels.

\subsection{Chain Security}

\begin{theorem}[Audit Chain Security]
Under SHA-256 (FIPS PUB 180-4~\citep{fips180}), an adversary cannot
produce a modified input file that has the same SHA-256 hash as the
original file without performing at least $2^{80}$ hash operations
(birthday bound for collision resistance).
At $10^{15}$ hash operations per second on current hardware, this
requires $>3 \times 10^7$ years of computation.
\end{theorem}

This is not a formal proof of SHA-256 collision resistance (which remains
an open problem) but reflects the current state of cryptographic knowledge.
No practical collision attack on SHA-256 is known as of 2026.

\subsection{Verification by Third Parties}

The \dia{} requires the redistricting manifest to be filed publicly
with the Secretary of State.
Any party---opposing counsel, expert witnesses, journalists, academic
researchers---can independently verify the manifest using:
\begin{verbatim}
bisect label-verify <label> --year <year> --manifest <path>
\end{verbatim}
This converts the audit chain from a self-certification into a
publicly verifiable claim.
